\documentclass[11pt,a4paper]{article}
\usepackage[utf8]{inputenc}
\usepackage{graphicx}
\usepackage{hyperref}
\usepackage{longtable}
\usepackage{booktabs}

\title{Software Design Document (SDD)}
\author{ADNS Project Team}
\date{November 25, 2025}

\begin{document}

\maketitle

\section{Introduction}

\subsection{Purpose}

This Software Design Document describes the software architecture and detailed design for the Autonomous Drone Navigation System (ADNS).

\subsection{Scope}

The ADNS software consists of:
\begin{itemize}
    \item Perception Module (sensor processing)
    \item Navigation Module (path planning, state estimation)
    \item Mission Control Module (autonomous mission execution)
    \item Communication Module (telemetry, commands)
\end{itemize}

\subsection{Definitions and Acronyms}

\begin{table}[h]
\centering
\begin{tabular}{ll}
\toprule
\textbf{Term} & \textbf{Definition} \\
\midrule
ROS2 & Robot Operating System 2 - middleware framework \\
DDS & Data Distribution Service - ROS2 transport layer \\
EKF & Extended Kalman Filter - state estimation algorithm \\
SLAM & Simultaneous Localization and Mapping \\
\bottomrule
\end{tabular}
\end{table}

\section{System Architecture}

\subsection{Architecture Overview}

The ADNS software follows a layered architecture:

\begin{enumerate}
    \item \textbf{Hardware Abstraction Layer}: Sensor drivers, actuator interfaces
    \item \textbf{Middleware Layer}: ROS2 communication, message passing
    \item \textbf{Application Layer}: Perception, navigation, mission control
    \item \textbf{User Interface Layer}: Ground station GUI, CLI tools
\end{enumerate}

\subsection{Module Decomposition}

\subsubsection{Perception Module}

\textbf{Responsibilities}:
\begin{itemize}
    \item Process LiDAR point clouds (100,000 points at 10 Hz)
    \item Process camera images (1920x1080 at 30 FPS)
    \item Fuse sensor data into unified world model
    \item Detect and classify obstacles
\end{itemize}

\textbf{Interfaces}:
\begin{itemize}
    \item Input: /lidar/points (sensor\_msgs/PointCloud2)
    \item Input: /camera/image (sensor\_msgs/Image)
    \item Output: /perception/obstacles (custom\_msgs/ObstacleArray)
\end{itemize}

\subsubsection{Navigation Module}

\textbf{Responsibilities}:
\begin{itemize}
    \item Estimate 6-DOF pose (position + orientation)
    \item Plan collision-free trajectories
    \item Execute waypoint navigation
    \item Handle dynamic obstacle avoidance
\end{itemize}

\textbf{Components}:
\begin{itemize}
    \item State Estimator (EKF)
    \item Path Planner (A*)
    \item Trajectory Generator
    \item Controller Interface
\end{itemize}

\section{Detailed Design}

\subsection{State Estimation (Extended Kalman Filter)}

The EKF estimates the drone state vector:

\[
\mathbf{x} = [p_x, p_y, p_z, v_x, v_y, v_z, q_w, q_x, q_y, q_z]^T
\]

Where:
\begin{itemize}
    \item $p_{x,y,z}$ = position (meters)
    \item $v_{x,y,z}$ = velocity (m/s)
    \item $q_{w,x,y,z}$ = orientation quaternion
\end{itemize}

\textbf{Prediction Step}:

\[
\hat{\mathbf{x}}_{k|k-1} = f(\mathbf{x}_{k-1}, \mathbf{u}_k)
\]

\[
\mathbf{P}_{k|k-1} = \mathbf{F}_k \mathbf{P}_{k-1} \mathbf{F}_k^T + \mathbf{Q}_k
\]

\textbf{Update Step}:

\[
\mathbf{K}_k = \mathbf{P}_{k|k-1} \mathbf{H}_k^T (\mathbf{H}_k \mathbf{P}_{k|k-1} \mathbf{H}_k^T + \mathbf{R}_k)^{-1}
\]

\[
\hat{\mathbf{x}}_k = \hat{\mathbf{x}}_{k|k-1} + \mathbf{K}_k (\mathbf{z}_k - h(\hat{\mathbf{x}}_{k|k-1}))
\]

\subsection{Path Planning (A* Algorithm)}

The A* algorithm finds the optimal path from start to goal:

\begin{enumerate}
    \item Initialize open set with start node
    \item While open set not empty:
    \begin{enumerate}
        \item Select node with lowest $f(n) = g(n) + h(n)$
        \item If goal reached, reconstruct path
        \item Expand neighbors, update costs
    \end{enumerate}
\end{enumerate}

\textbf{Heuristic Function}:

\[
h(n) = \sqrt{(x_n - x_{goal})^2 + (y_n - y_{goal})^2 + (z_n - z_{goal})^2}
\]

\section{Data Structures}

\subsection{Obstacle Representation}

\begin{verbatim}
struct Obstacle {
    uint32_t id;
    geometry_msgs::Point position;
    geometry_msgs::Vector3 velocity;
    geometry_msgs::Vector3 dimensions;
    enum Type { STATIC, DYNAMIC } type;
    float confidence;
};
\end{verbatim}

\subsection{Waypoint}

\begin{verbatim}
struct Waypoint {
    geometry_msgs::Point position;
    float tolerance;  // meters
    float speed;      // m/s
    bool hover;       // hover at waypoint?
    float hover_time; // seconds
};
\end{verbatim}

\section{Performance Requirements}

\begin{table}[h]
\centering
\begin{tabular}{lll}
\toprule
\textbf{Module} & \textbf{Requirement} & \textbf{Target} \\
\midrule
Perception & Point cloud processing & < 50 ms \\
Perception & Image processing & < 30 ms \\
Navigation & State estimation rate & 100 Hz \\
Navigation & Path planning & < 200 ms \\
Mission Control & Command response & < 100 ms \\
\bottomrule
\end{tabular}
\end{table}

\section{Error Handling}

\subsection{Sensor Failures}

\begin{itemize}
    \item \textbf{LiDAR failure}: Switch to camera-only mode, reduce speed
    \item \textbf{Camera failure}: Rely on LiDAR, disable visual odometry
    \item \textbf{IMU failure}: Trigger emergency landing
\end{itemize}

\subsection{Communication Failures}

\begin{itemize}
    \item \textbf{Comms loss $>$ 10 seconds}: Execute return-to-home
    \item \textbf{Comms latency $>$ 500 ms}: Alert operator, continue mission
\end{itemize}

\section{Safety-Critical Functions}

The following functions are classified as DO-178C Level C:

\begin{itemize}
    \item Obstacle detection and avoidance
    \item Geofence enforcement
    \item Emergency landing controller
    \item Watchdog monitoring
\end{itemize}

\textbf{Verification}: MC/DC (Modified Condition/Decision Coverage) required for all safety-critical code.

\section{Testing Strategy}

\subsection{Unit Tests}

\begin{itemize}
    \item Google Test framework
    \item Target: 95\% statement coverage
    \item Target: 85\% MC/DC coverage (safety-critical code)
\end{itemize}

\subsection{Integration Tests}

\begin{itemize}
    \item ROS2 integration testing (launch tests)
    \item Hardware-in-the-Loop (HIL) testing
\end{itemize}

\section{Document Control}

\textbf{File}: 21\_SDD\_ADNS\_Software.tex \\
\textbf{Version}: 1.0.0 \\
\textbf{Date}: November 25, 2025 \\
\textbf{Status}: Approved \\
\textbf{Classification}: Company Confidential

\end{document}
